\section{Discussion}\label{sec:discussion}

\subsection{Main Findings}

This study asked whether digital factor allocation promotes regional coordinated development and whether the relationship is nonlinear. The answer is yes on both counts, with qualifications. On average DFA is positively associated with RCD, and the estimate of 0.241 survives alternative weighting of the outcome, the exclusion of municipalities and of 2020, and winsorization. The association is stronger in provinces that have crossed a critical level of economic development and is concentrated in the west. A financial-development threshold indicates declining returns at very high levels of financial depth, although that result rests on a small group of municipalities.

Three points follow for the literature. First, the digital dividend is conditional: digital inputs are more productive when they can be combined with skills, markets, and institutions. Second, the digital-divide and digital-dividend views are both right at different stages of development: access without absorptive capacity does not convert into coordination gains, and once the capacity is in place the same inputs do more. Third, a composite index conceals divergent behaviour across its components. Digital finance carries most of the association, while industrial digitalization and digital human capital change sign across the threshold.

The spatial analysis narrows the claim. The positive association holds in every specification and in both regimes, and the threshold location is stable when the outcome is allowed to depend on neighbouring outcomes. The size of the gap between regimes is not stable. The threshold therefore describes a development gradient in the cross-section, not a sharp structural boundary.

\subsection{Absorptive Capacity and the Digital Divide}

The economic-development threshold is consistent with absorptive capacity. Digital technologies are general-purpose technologies whose effects depend on organizational change, human capital, and complementary infrastructure \citep{brynjolfsson2014,omahony2005}. In low-income provinces digital infrastructure may be used mainly for communication and consumption, while productive and coordination effects require inputs that are in short supply. In higher-income provinces digital inputs can be combined with skilled labour, financial services, and market institutions.

This is also consistent with the multi-level digital divide. Access is the first level, skills the second, and the ability to convert access into material benefits the third \citep{vandijk2005,norris2001,warschauer2003}. Our results suggest that the third-level divide binds in low-income provinces, so providing access alone is not enough.

The regional and threshold results are not in conflict. The regional estimate compares average outcomes across regions over the sample period, while the threshold estimate compares provinces at different development stages within a period. The west can have a larger average association because it starts from a lower base even if the marginal association at a given development level is smaller.

\subsection{Financial Development and Diminishing Returns}

The financial-development threshold suggests that the marginal association of DFA declines at very high levels of financial depth, consistent with diminishing returns to digital finance at the extensive margin. Where financial systems are shallow, digital finance expands access to basic services; where they are deep, it adds less to access and may instead intensify competition or encourage speculative activity \citep{beck2007}. The result rests on 30 high-regime observations, mostly Beijing, Shanghai, and Tianjin, and should be read as a caution against assuming that more digital finance always produces more coordination.

The decomposition of the index points at the same factor from another direction. Entered jointly, digital finance is the only dimension with an independent association with RCD, and its regime pattern mirrors the composite: positive in both regimes and slightly larger above the threshold. Financial depth therefore matters twice, as a moderator of the association and as the carrier of most of it. The two roles are conceptually distinct but point to the same policy margin: the quality and direction of digital finance rather than its volume.

\subsection{Regional Heterogeneity}

The association is strongest in the west and weaker or insignificant elsewhere under two-way fixed effects, which supports the inclusive-digitalization view: where digital infrastructure and financial access are scarce, additional digital factors have larger marginal effects. The northeast is weaker and sometimes negative. It has relatively high initial urbanization and public services but a slow digital transition and a declining working-age population, so digital factor allocation may not offset these structural constraints without a broader revitalization strategy.

\subsection{Policy Implications}

Digital factor allocation should be sequenced with economic development. In low-income provinces the priority is basic connectivity, digital literacy, and complementary public services; advanced digital industry and digital finance are unlikely to raise coordination until those foundations are in place. In higher-income provinces the priority shifts to deeper integration of digital technologies into production, public services, and urban-rural linkages.

Policies should be tailored to regional absorptive capacity. The west has the strongest average association and should be a priority for digital infrastructure and digital finance investment, paired with human capital development and market reform. Because the short-run association with the urban-rural income ratio is positive, gains may be captured disproportionately by urban and skilled groups; rural e-commerce, digital agriculture, digital public services, and targeted skills training can broaden them. Advanced regions should attend to the quality and direction of digital investment rather than its quantity, given the declining association at high financial depth.

\subsection{Limitations}\label{sec:limitations}

Four limitations deserve emphasis. First, the analysis covers 30 provinces and does not capture within-province heterogeneity; city-level data would provide a more granular picture. Second, the indices depend on indicator selection and weighting. We use alternative weighting schemes and alternative outcome definitions, but the indices remain simplifications.

Third, the estimates are conditional associations. DFA may be endogenous to RCD, and the lagged specification is not significant, which suggests a contemporaneous relationship that could reflect reverse causality. Oster bounds and the randomisation placebo indicate that unobservables would have to be more important than the observed controls to eliminate the association, but they do not identify a causal effect. External instruments such as historical telephone penetration, terrain ruggedness, or distance to the backbone network are not available for this sample, and we do not construct a pseudo-instrument from the index components. Future work should exploit such variation.

Fourth, the threshold is sensitive to specification. Region-specific trends, year fixed effects, and spatial specifications either reduce its precision or move its location, and the size of the break is not stable across spatial models. The threshold should be read as a long-run pattern, not a short-run causal effect. The sub-dimension slopes are also imprecise because the low-income regime contains 66 observations, and the income series carry measurement error from the household survey and from the urban-rural classification. Finally, the sample includes the COVID-19 period; the results are robust to dropping 2020, but 2021 may still be affected.
